\documentclass[letterpaper]{article} % DO NOT CHANGE THIS
\usepackage[submission]{aaai2027}  % DO NOT CHANGE THIS
% The serif, sans-serif, and monospaced fonts are loaded automatically by
% aaai2027.sty (newtxtext, helvet, courier). DO NOT add \usepackage{times},
% \usepackage{helvet}, \usepackage{courier}, or any other font package.
\usepackage[hyphens]{url}  % DO NOT CHANGE THIS
\usepackage{graphicx} % DO NOT CHANGE THIS
\urlstyle{rm} % DO NOT CHANGE THIS
\def\UrlFont{\rm}  % DO NOT CHANGE THIS
\usepackage{natbib}  % DO NOT CHANGE THIS AND DO NOT ADD ANY OPTIONS TO IT
\usepackage{caption} % DO NOT CHANGE THIS AND DO NOT ADD ANY OPTIONS TO IT
\frenchspacing  % DO NOT CHANGE THIS
\usepackage{booktabs}
\usepackage{amsmath}
\usepackage{amssymb}
\usepackage{multirow}


\usepackage{xcolor}
\usepackage{tcolorbox}
\tcbuselibrary{skins}

\newtcolorbox{findingbox}[1][]{
  colback=gray!5,
  colframe=gray!40,
  leftrule=3pt,
  rightrule=0pt,
  toprule=0pt,
  bottomrule=0pt,
  arc=0pt,
  boxsep=2pt,
  left=6pt,
  right=6pt,
  top=4pt,
  bottom=4pt,
  #1
}

% Placeholder for cross-simulator numbers that are filled in from crosssim/results (remove when final).
\newcommand{\xs}[1]{\textcolor{red}{#1}}

\pdfinfo{
/TemplateVersion (2027.1)
}

\setcounter{secnumdepth}{2}

\title{The \textit{Silicon Society} Cookbook: A Controlled Study of the Design Space of LLM-Based Social Simulations}

\author{
    Anonymous Submission
}
\affiliations{
    ~
}

\begin{document}

\maketitle

\begin{abstract}

\textit{If you know your model and know yourself, you need not fear the result of a hundred experiments.}
Numerous recent studies simulate human social behavior using LLM-based \textit{Silicon Societies}, yet the impact of the design decisions on the outcomes of these society models remains poorly characterized. Under these circumstances, artifacts stemming from arbitrary design choices in seemingly irrelevant system components may be incorrectly attributed to genuine social effects. Given that some of these models aim to inform future policy-making decisions, these misattributions could be of consequence. To answer these concerns, we study the internal sensitivity of one such model in depth, then test whether its main findings carry over to three independently built simulators (OASIS, Concordia and SiliSocS). We characterize parts of the \textit{Silicon Society} design space with a large sweep of 1{,}224 simulation roll-outs across five experiments while systematically varying 10 design parameters. We find that the single largest driver of survey-response dynamics is whether a model is \emph{fine-tuned on social media data} or not (partial $\eta^2$ up to $0.89$). The space is partially non-additive: notably, several axes act only through model-gated interactions. Through several null model baselines, we decompose which parts of the measured effects are socially driven, revealing that a large share of measured answer change is prompt-perturbation noise or a mechanical by-product of how the metrics are computed. Because our agents are stateless, we treat these as changes in survey responses, not as belief revision. We believe that adoption of our framework for robustness checks could strengthen the field of LLM-based social sciences.
\end{abstract}
\section{Introduction}
Being wrong is bad, but getting to a correct answer through incorrect means may sometimes be worse. Agent-based models have a long history in the social sciences \citep{de2014agent}, but in recent years a new paradigm  has emerged where Large Language Models (LLMs) are used to simulate human interaction. The promise of systems that capture both micro-level behavior (users, messages) and macro-level phenomena (echo chambers, herd effects, diffusion) at arbitrary scale has proven alluring. New and increasingly complex simulators are published continuously, yet a chronic lack of validation limits the field's ability to build cumulatively and to address the known weaknesses of classical agent-based models \citep{ye2026stopdrawingscientificclaims,li2025llmgeneratedpersonapromise,larooij2025largelanguagemodelssolve,puelmatouzel2026validationgap}. Demonstrating that these simulators faithfully reproduce human behavior and that their findings generalize beyond their original setting remains an open challenge. Yet even this may not suffice: LLM-based social simulations comprise numerous explicit and implicit design choices whose influence on outcomes is poorly understood. Apparent alignment with human social dynamics may therefore stem from unintended interactions between design choices rather than the mechanisms the simulator is meant to capture, producing convincing but misleading results and unwarranted confidence in a model's predictions \cite{ye2026stopdrawingscientificclaims}. Progress therefore requires not only validating these systems against human data, but also understanding the effects of individual design decisions, how those decisions interact, and developing standardized methods to identify and mitigate false positives.


We study \textit{Silicon Societies}, defined as simulations of human interactions that use LLMs as agents \cite{puelmatouzel2026validationgap}. Examples include controlled research simulators such as OASIS \citep{yang2025oasisopenagentsocial}, AgentSociety \cite{piao2026agentsocietylargescalesimulationllmdriven} and Social Simulacra \cite{park2022socialsimulacracreatingpopulated}. We exclude multi-agent frameworks whose goal is task completion rather than human-likeness. While some of our conclusions may extend to naturally emerging LLM-populated platforms, which became relevant following Moltbook's meteoric rise to fame \cite{brach2026moltbookfilesharmlessslopocalypse}, they are not our main focus.


Our study spans 5 experiments over 1,224 simulation roll-outs in our own simulator, plus a cross-simulator replication in OASIS, Concordia and SiliSocS. Our measured quantity is the answer agents give to a periodically administered survey question. Because the agents have no persistent state beyond their context window, a change in answer reflects a change in what is in the context (or in how the question is phrased), not the revision of a held belief; we use the term \emph{survey-response dynamics} for this reason. We hypothesize that the design space is non-additive ($X_{A+B} \neq X_A + X_B$), exhibiting complex parameter interactions that make outcomes hard to predict from initial settings alone.


Our findings reveal that the design space is neither uniformly additive nor uniformly complex: In particular, several effects appear to be model-specific. Null models are central to our interpretive framework. A scrambled-stimulus baseline reveals that much of what raw survey-response metrics register as social influence is in fact prompt-perturbation noise: agents change their survey answers after reading shuffled, meaningless content nearly as often as after genuine social stimuli. A separate response-shuffling null shows that, in our setup, the network's apparent influence on neighbor alignment is largely mechanical, a consequence of graph structure in how the metric is computed rather than of agents' answers tracking their neighbors'. This is a statement about what these metrics can detect in the regime we study, not a claim that topology cannot shape LLM social simulations. In both cases a naive reading of the raw measurements would overstate social effects; the null models are what let us say which part of the signal is socially driven.

%Assessing whether such systems are \emph{valid} models of human behavior is one of the central open problems of the field, and precisely because it is unsolved we make no such claim here. Instead we ask a prior and more tractable question: \textbf{when we build a Silicon Society, which design choices actually change what happens inside it, and how predictably?} This matters for two reasons. First, the number of LLM social-simulation papers is exploding: a Google Scholar search for ``LLM social simulation'' returns thousands of results for 2026 alone, and these systems are increasingly proposed for practical, decision-relevant use. If decisions come to be informed by such simulations, it is essential that the simulations be well understood. Second, several works explicitly call for exactly this kind of foundational scrutiny of design and measurement before realism claims can be trusted \citep{larooij2025largelanguagemodelssolve,li2026positionaiagentsyet,seshadri2026lost}.

%\paragraph{Scope.} We characterize the \emph{internal sensitivity} of a controlled simulator: how its measured outcomes respond to design choices. We do not claim behavioral or societal human-likeness, and we avoid ``realism'' language throughout. Because LLM agents in our simulator (and in most comparable ones) carry no persistent belief state beyond their context window and adapter weights, we deliberately use the term \emph{survey-response dynamics} rather than ``opinion dynamics'': what we measure is how a stateless LLM population's survey answers respond to stimuli, context, and structure, which is a well-defined and scientifically informative object even though it is not belief revision. All empirical conclusions are scoped to controlled simulators of the architecture we study; the broad ``Silicon Society'' framing motivates the problem but does not extend the claims.

We summarize our contributions as follows:
% \paragraph{Contributions.}
\begin{itemize}
\item \textbf{Validity Recipe:}  We present reproducible internal sensitivity and measurement-validity instrumentation in four components: \textbf{parameter sensitivity} experiments, \textbf{parameter interaction} experiments, \textbf{null models} and rigorous \textbf{statistical analysis} of outcomes.
\item \textbf{Impactful Design Choices:} By applying our methods, we identify that the design choices with the largest effects on simulation outcomes and dynamics are \textbf{finetuning on social-media data} and the \textbf{number of simulated agents}, and quantify the direction and magnitude of these effects.
\item \textbf{Complex Interactions:} We show that the simulation design space is \textbf{partially non-additive}.
\item \textbf{Behavior or Artifact?:}  Using null models, we distinguish socially driven response changes from effects attributable to noise or simple mechanistic explanations.
\item \textbf{Beyond One Simulator:} We re-run the central tests in three independently built simulators (OASIS, Concordia, SiliSocS) with the same base models and adapters, and report which findings transfer and which are tied to our implementation.
\end{itemize}

We release our code and experimental pipeline to support the reproducibility of our methods.

\section{Related Work}


\paragraph{LLM-based social simulators.} Works in this area are numerous; \citet{larooij2025largelanguagemodelssolve} survey dozens of simulators with human-imitation objectives. Notable systems include OASIS \citep{yang2025oasisopenagentsocial}, which scales to one million agents and reports macro-scale effects such as herding and echo chambers, and Concordia \citep{vezhnevets2023generative}, which introduced modular cognitive agent components. Other recent work follows varied settings and validation strategies \citep{zheng2026rolesimllm,lu2026real,yuzhe2026twinmarket}. In parallel, LLM-populated networks have begun to appear outside controlled settings and have also attracted study \citep{sodano2026emergencefragilityllmbasedsocial,jiang2026humans,holtz2026anatomy}.

\paragraph{The Validation Gap.} The mismatch between how fast simulators are published and how weakly they are validated is increasingly noted. \citet{larooij2025largelanguagemodelssolve} argue that poor standards prevent LLM social simulation from addressing long-standing agent-based-model weaknesses; \citet{li2025llmgeneratedpersonapromise} call for a science of persona creation; \citet{li2026positionaiagentsyet} argue that current validation practice does not meet what simulation-as-a-science requires; and \citet{seshadri2026lost} warn that evaluation practices risk misrepresenting agent capabilities. Our study addresses a prerequisite of validation: understanding which design decisions shape simulator behavior. \citet{ye2026stopdrawingscientificclaims} notably argue for robustness audits before drawing scientific conclusions and test the impact of perturbations, but at a much smaller scale than ours, without measuring the relative importance of design choices or quantifying their interactions. Our study is complementary to their demonstration of the problem's relevance.

\paragraph{Network Topology and Opinion Dynamics.} Outside LLM simulation, the effect of network structure on opinion dynamics has been studied extensively in computational social science, network science, and statistical physics. Classical models (voter, Deffuant, Hegselmann--Krause, DeGroot/Friedkin--Johnsen) consistently show that structural properties such as degree heterogeneity, clustering, community structure, and path length strongly influence consensus, polarization, fragmentation, and diffusion. We do not aim to rediscover these results, but to ask whether the same structural dependencies re-emerge when agents communicate in natural language through an LLM rather than hand-designed update rules, and how large they are relative to other design choices.

\section{Simulator and Experimental Design}


%\subsection{Simulator}
Our simulator instantiates a set of LLM-backed agents and connects them with a directed followership network. Each agent is assigned a fixed activity probability drawn from a Zipfian distribution, so that a minority of ``power users'' account for most activity. At each step, ten agents are selected by activity probability; nine observe a recent thread in which someone they follow participated, and one either starts a new thread (probability $\frac{1}{3}$) or replies to an existing one (probability $\frac{2}{3}$). Runs last 2{,}500 steps with agents surveyed every 250 steps, yielding 11 survey snapshots. Every run is fully seeded from its design row. The appendix specifies decoding and generation settings, per-run seeding, the computing infrastructure (hardware and software versions), and example threads.

%We expose several mechanistic axes at negligible additional compute cost:
%\begin{itemize}
%\item \textbf{Activity exponent} $\in \{0.0, 0.5, 1.0\}$: the Zipf exponent of the activity distribution, i.e. an opinion-leader / power-user axis ($0$ is egalitarian, $1$ is strongly leader-driven).
%\item \textbf{Number of news agents} $\in \{0, 1, 4\}$: a biased broadcaster that posts pre-generated messages favoring one survey option, placed at the highest-degree node(s); a dose-response generalization of a single-broadcaster condition.
%\item \textbf{Survey cadence} (\texttt{max\_steps}, \texttt{survey\_interval}): to test per-agent-matched surveying at different population sizes.
%\item \textbf{Stimulus mode} $\in \{\text{normal}, \text{scrambled}, \text{none}\}$: measurement-validity baselines (Section~\ref{sec:measurement}).
%\item \textbf{Dual-order surveying}: every survey is asked in canonical and option-order-reversed phrasings, recording both answers and their log-probability margins.
%\end{itemize}
%Every run is fully seeded (Python/NumPy/PyTorch) from its design row, and the empirical mean degree of the generated graph is logged per run.

\subsection{Fine-tuning on Social-Media Data}
To move agents away from the characteristic ``assistant'' register and toward social-media discourse, we fine-tune base models on the BluePrint dataset \citep{bückkaeffer2025textttblueprintsocialmediauser} using Low-Rank Adaptation \citep{hu2021loralowrankadaptationlarge}. BluePrint clusters users into 25 behavioral archetypes; we train one adapter per archetype for each base model, so that a population can be composed from adapters in chosen proportions. We use four base models: Llama-3.1-Minitaur-8B \citep{binz2024centaurfoundationmodelhuman}, Llama-3.1-8B (to isolate the effect of Minitaur's post-training), gemma-3-4b-pt \citep{gemma_2025}, and Qwen2.5-7B-Instruct \citep{qwen2}. We measure alignment with human opinion using SimBench \citep{hu2025simbenchbenchmarkingabilitylarge} before and after fine-tuning; fine-tuning improves SimBench scores on average except where the base was already strong. The appendix details the dataset, the LoRA training setup, the hyperparameter ranges we swept and how final settings were chosen, and per-model SimBench scores.

\subsection{Balanced Design}
\label{sec:design}
%All runs are rows of committed design files generated deterministically from a single seed. 
Marginals and model$\times$parameter cells are balanced. In particular, we cross a no-adapter (base-only) arm across all four base models. Table~\ref{tab:experiments} lists the five experiments.

\begin{table*}[htbp!]
\centering
\small
\setlength{\tabcolsep}{5pt}
\begin{tabular}{llp{10cm}rr}
\toprule
ID & Purpose & Parameters varied & Runs & Agents \\
\midrule
E1 & Core balanced sweep & Base model, proportions, question, population size, graph family, homophily, survey-in-context, news dose, activity exponent (fine-tuning held on) & 576 & 64 -- 1,024 \\
E2 & Cross fine-tuning with model & Fine-tuning (on / off) $\times$ base model & 144 & 256 \\
E3 & Controlled topology study & Graph family (7) $\times$ base model $\times$ question & 336 & 256 \\
E4 & Measurement baselines & Stimulus mode (normal / scrambled / none) $\times$ base model $\times$ survey-in-context & 96 & 256 \\
E5 & Survey-cadence control & Population size, with the survey interval matched to $N$ & 72 & 64 -- 1,024 \\
\midrule
\multicolumn{3}{l}{\textbf{Total}} & \textbf{1{,}224} & \\
\bottomrule
\end{tabular}
\caption{The five experiments and the parameters each varies.}
\label{tab:experiments}
\end{table*}

% MOVED TO APPENDIX (Design & Power note): minimum-detectable-effect calculation.
% With 576 balanced core runs, a two-level parameter has 288 runs per level, giving a minimum detectable effect of $d \approx 0.38$ at $\alpha = 0.001$ with power $0.9$; four-level parameters have 144 per level ($d \approx 0.46$ for pairwise contrasts). Every two-way interaction cell holds at least 36 runs, and the replicated E3/E4 designs provide at least 8 runs per primary cell.

\subsection{Parameters}
We vary ten parameters. Eight are self-contained:
\begin{itemize}
\setlength{\itemsep}{1pt}
\item \textbf{Base model} (4 models): the underlying LLM generating every agent's messages and survey answers.
\item \textbf{Fine-tuning} (on / off): whether each base model wears its BluePrint LoRA adapter, fine-tuned on social-media data, or runs unmodified.
\item \textbf{Question} (Q1, Q2, Q3): which of three survey questions the agents are polled on throughout the run.
\item \textbf{Population size $N$} (64, 256, 1024): the number of agents.
\item \textbf{Homophily} (on / off): whether agents giving the same initial survey answer are preferentially connected at initialization, seeding clusters of like-answering agents.
\item \textbf{Survey-in-context} (on / off): whether an agent's own previous survey answer is placed in its context window.
\item \textbf{News dose} (0, 1, 4): the number of broadcaster agents injecting a shared news message into the network.
\item \textbf{Activity exponent} (0.0, 0.5, 1.0): the skew of the per-agent posting-frequency (Zipf) distribution; $0$ is egalitarian, $1$ is strongly opinion-leader-driven.
\end{itemize}
The remaining two, \emph{proportions} and \emph{graph family}, require more details.

\paragraph{Proportions.} This parameter sets how the 25 archetype adapters are mixed to form the agent population. \emph{Uniform} assigns each adapter exactly $1/25$ of the population, and \emph{blueprint} reproduces the empirical archetype frequencies observed in the BluePrint dataset. \emph{Average} and \emph{distribution} are opinion-targeted schemes. Each of the 25 archetype adapters induces, for every survey question, a probability vector over that question's answer options; stacking these vectors across all questions gives an answer matrix $M_a$ for adapter $a$. Human responses, stacked the same way, give a target matrix $H$. A population with adapter mixing weights $w \in \Delta^{24}$ (a distribution over the 25 adapters) then has expected answer matrix $\sum_a w_a M_a$, and we pick $w$ to minimize its distance to the human target,
\begin{equation}
  \min_{w \in \Delta^{24}} \; \Big\| \textstyle\sum_a w_a M_a - H \Big\|_2^2,
\end{equation}
a convex program in $w$. For \emph{average}, $H$ stacks the modal human answer per question; for \emph{distribution}, it stacks the full empirical human answer distribution. We revisit the Average-vs-Distribution contrast in Section~\ref{sec:proportions}.

\paragraph{Graph family.} We test seven graph families (E3): Erd\H{o}s--R\'enyi, Barab\'asi--Albert, stochastic block, cycle, power-law cluster, fully connected, and forest fire, crossed with four base models and three questions, with the five parameterizable families density-matched to mean degree $\approx 16$. We test each metric with a permutation $\eta^2$ that shuffles topology labels within (model $\times$ question) strata.

\subsection{Metrics and Terminology}
\label{sec:metrics}
All metrics are computed from model survey responses; they carry no human baseline and therefore measure internal sensitivity, not realism. Our four headline behavioral metrics are defined in Table~\ref{tab:metrics}.

\paragraph{What a survey answer is (and is not).} Agents carry no persistent belief state: between surveys, all that changes is the text in their context window. An agent's survey answer at a snapshot is therefore the option the model scores highest \emph{given its current context}, and a change in answer means the context changed enough to move that choice. It is not evidence that a belief was revised. We call the object of study \emph{survey-response dynamics}, and we use ``shift'', ``alignment'' and ``consensus'' strictly as descriptions of answer patterns across snapshots. Where we use ``opinion'' in a metric name inherited from the literature, it should be read in this sense.

\begin{table*}[t]
\centering
\small
\begin{tabular}{@{}llp{10cm}c@{}}
\toprule
Acr. & Full name & Meaning & Formula \\
\midrule
OSR & Option Shift Rate & Fraction of agents whose survey answer differs from their previous answer. & $N_{\text{chg}}/N_{\text{total}}$ \\[4pt]
MFR & Majority Follow Rate & Among agents who changed their answer, the fraction that joined the majority. & $N_{\text{chg}\to\text{maj}}/N_{\text{chg}}$ \\[4pt]
NASR & Neighbor Alignment Shift Rate & Share of agents who switched to their neighbors' majority answer. & $N_{\text{chg}\to\text{nb}}/N_{\text{total}}$ \\[4pt]
NCC & Net Consensus Change & Change in the population's modal-answer share from the first to the last snapshot; positive = the population converged, negative = it fragmented. & $C_{t_f} - C_{t_0}$ \\
\bottomrule
\end{tabular}
\caption{The four headline behavioral metrics. $C = N_{\max}/N_{\text{total}}$ is the population consensus (modal-answer share) at a snapshot; $N_{\text{chg}}$ is the number of agents who changed answer, of which $N_{\text{chg}\to\text{maj}}$ moved toward the global majority and $N_{\text{chg}\to\text{nb}}$ toward their neighbor majority.}
\label{tab:metrics}
\end{table*}

Alongside these we track structural metrics: network assortativity, mean local agreement, and cross-cutting edge fraction, together with their change from the first to the last snapshot (e.g.\ $\Delta$ assortativity). From the dual-order surveys we additionally compute an \emph{order-consistency rate} (the fraction of answers that survive reversing the option order), the mean absolute log-probability margin, and the low-margin flip share. Exact definitions of all metrics are in the appendix.

\subsection{Statistical Protocol}
\label{sec:stats}
We attribute variance with a Type-III ANOVA using sum-to-zero contrasts over all main effects and all two-way interactions, and report \emph{partial} $\eta^2$ rather than single-parameter $\eta^2$ (which would let a parameter absorb correlated choices). Because the base-only condition perfectly aliases the proportions parameter, we run two passes: a \emph{pooled} pass (E1$+$E2, 720 runs) in which fine-tuning status is crossed with model and proportions is dropped, and an \emph{E1-only} pass (576 runs) in which proportions is identifiable and fine-tuning is held constant. Each headline states its pass. Within every metric family we report Benjamini--Hochberg FDR $q$-values, bootstrap percentile confidence intervals on each partial $\eta^2$, and a mixed-effects refit with the compute batch as a random intercept so that batch-correlated noise cannot masquerade as a design effect. We report $q$-values rather than relying on a threshold.


\section{Results}


For each result we report the effect size (partial $\eta^2$ with bootstrap CI and FDR $q$). Where a finding does not appear to have a mechanistic explanation and survives robustness checks, we flag it as more likely to generalize; where it plausibly depends on our simulator's design we flag it as implementation-tied.

\subsection{Which Parameters Move the Needle?}
\label{sec:attribution}

\begin{figure*}[t]
  \centering
  \includegraphics[width=\textwidth]{Figures/fig_eta2_dotplot.pdf}
  \caption{Main-effect partial $\eta^2$ (with 95\% bootstrap CIs) for each design parameter across the four outcome metrics. Fine-tuning and number of agents $N$ dominate; topology, homophily, and activity exponent are negligible in the range of setups we explore.}
  \label{fig:dotplot}
\end{figure*}

Table~\ref{tab:attribution} reports the largest terms per metric (including interactions), and Figure~\ref{fig:dotplot} places every parameter's main effect side by side across the four metrics. Both tell the same story: the dominant driver of survey-response dynamics is \emph{fine-tuning status}, followed by the number of agents. Turning the BluePrint adapters on or off explains the majority of the variance in every response-intensity metric: partial $\eta^2 = 0.75$ for OSR, $0.89$ for MFR, and $0.78$ for NASR, each with $q < 10^{-180}$ and tight bootstrap intervals. The graph-type main effect, by contrast, is small (partial $\eta^2 \leq 0.001$ for all metrics, not significant). Within the range we swept (density-matched graphs of up to 1{,}024 stateless agents over 2{,}500 steps, measured through survey answers), topology has no measurable impact on these metrics. This does not rule out topology effects in other regimes, such as much larger networks, agents with persistent memory, longer horizons, or outcomes other than survey answers (Section~\ref{sec:topology}).

Figure~\ref{fig:finetune_direction} unpacks the effect by sign, plotting the signed delta (mean with the LoRA adapter minus without) per base model on the $N=256$ region where both arms overlap. Fine-tuning raises the response-intensity metrics for the Llama models but reverses sign for Qwen: the effect is not a uniform ``fine-tuning makes agents more reactive'' shift but a model-gated one, whose direction cannot be read off the main effect alone.

\begin{table}[t]
\centering
\small
\setlength{\tabcolsep}{3.5pt}
\begin{tabular}{lllc}
\toprule
Metric & Top terms & partial $\eta^2$ [95\% CI] & $q$ \\
\midrule
OSR & fine-tuning & 0.754 [0.72, 0.78] & $<10^{-180}$ \\
    & model$\times$question & 0.342 [0.27, 0.42] & $<10^{-50}$ \\
    & population & 0.340 [0.28, 0.40] & $<10^{-53}$ \\
\midrule
MFR & fine-tuning & 0.893 [0.88, 0.91] & $<10^{-290}$ \\
    & population & 0.589 [0.53, 0.65] & $<10^{-115}$ \\
    & fine-tuning$\times$pop. & 0.579 [0.52, 0.65] & $<10^{-113}$ \\
\midrule
NASR & fine-tuning & 0.783 [0.75, 0.81] & $<10^{-200}$ \\
    & population & 0.339 [0.27, 0.40] & $<10^{-53}$ \\
    & model$\times$question & 0.328 [0.26, 0.40] & $<10^{-47}$ \\
\midrule
NCC & model$\times$question & 0.247 [0.20, 0.31] & $<10^{-32}$ \\
    & model$\times$fine-tuning & 0.060 [0.03, 0.11] & $<10^{-6}$ \\
    & model (main) & 0.006 [0.00, 0.03] & 0.64 (ns) \\
\bottomrule
\end{tabular}
\caption{Top 3 most impactful terms per metric (largest Type-III partial-$\eta^2$ terms, pooled pass E1$+$E2, 720 runs). A$\times$B denotes an interaction between A and B. Fine-tuning status dominates response-intensity metrics. Full tables in the appendix.}
\label{tab:attribution}
\end{table}

\begin{figure*}[t]
  \centering
  \includegraphics[width=\textwidth]{Figures/fig_finetune_direction.pdf}
  \caption{Fine-tuning on social-media data intensifies survey-response dynamics. Signed delta (mean with LoRA adapter $-$ mean without) per metric and base model, restricted to the $N=256$ design region where both arms overlap (95\% bootstrap CIs). Fine-tuning raises response-intensity metrics for the Llama models but reverses sign for Qwen, illustrating the model$\times$fine-tuning interaction.}
  \label{fig:finetune_direction}
\end{figure*}

Net Consensus Change tells the same story from the other direction: it is governed by a model$\times$question interaction (partial $\eta^2 = 0.247$), while the model main effect is not significant after FDR correction.
% MOVED TO APPENDIX: E1-only re-estimation of NCC. In the E1-only pass, where proportions becomes identifiable, the model$\times$question interaction again dominates NCC ($\eta^2_{\text{partial}} = 0.369$), followed by the model main effect ($0.205$) and a model$\times$proportions interaction ($0.144$).

Fine-tuning on versus off is the single largest lever on how much a population's answers move; base-model family matters, but chiefly through how it interacts with the question and the population composition rather than as a standalone effect, so we treat it as descriptive rather than causal-explanatory (its bundled confounds are enumerated in the limitations). The mixed-effects refit with compute batch as a random intercept leaves the fine-tuning and interaction effects intact, so they are not an artifact of batch structure. \emph{Flag: expected to generalize.}

\begin{figure}[t]
  \centering
  \includegraphics[width=\columnwidth]{Figures/fig_interaction_hist.pdf}
  \caption{A few interactions dominate. Distribution of two-way interaction partial $\eta^2$ (the off-diagonal cells of the per-metric parameter pair matrices, pooled across the four headline metrics). The mass is concentrated near zero (median $0.008$); only $20$ of the $176$ interaction terms clear $\eta^2 = 0.05$, and a short tail (model$\times$question, fine-tuning$\times$population, model$\times$proportions) carries the rest.}
  \label{fig:interaction_hist}
\end{figure}

Only a handful of parameter interactions actually matter. Figure~\ref{fig:interaction_hist} pools every two-way interaction partial $\eta^2$ across the four metrics; the distribution is sharply long-tailed (median $0.008$, only $20$ of $176$ terms above $0.05$). The design space is not a dense web of couplings: most parameter pairs are effectively independent, and the non-additivity that exists is carried by a small, nameable set (model$\times$question, fine-tuning$\times$population, model$\times$proportions). Full per-metric heatmaps appear in the appendix.

\subsection{Population Scale and Survey-Cadence Artifact}
\label{sec:scale}

Population size is a large term for structural metrics (partial $\eta^2 = 0.89$ for mean local agreement, with a comparably large fine-tuning$\times$population interaction), but part of this is mechanical: smaller populations have denser labeled neighborhoods. It is also sizable for OSR and MFR, which raises a worry that it reflects the fixed survey cadence: with more agents and a fixed 250-step interval, a smaller fraction of the population acts between snapshots, mechanically depressing measured shift rates.

E5 tests this directly by matching cadence to population (survey interval proportional to $N$). The OSR-versus-$\log_2 N$ slope flattens from $-0.0089$ per doubling under fixed cadence (Pearson $r = -0.16$, $n = 288$) to $-0.0022$ under matched cadence ($r = -0.03$, $n = 72$), a roughly fourfold reduction that removes the trend. The scale dependence of OSR is therefore largely a measurement artifact of survey cadence, not a property of the dynamics. \emph{Flag: implementation-tied (survey cadence).}

\subsection{Network Topology}
\label{sec:topology}

Pooled across models and questions, topology has a small but detectable association with neighbor alignment (NASR $\eta^2 = 0.048$, permutation $p = 2\times10^{-4}$) and majority-following (MFR $\eta^2 = 0.018$, $p = 2\times10^{-4}$), and no significant effect on NCC, OSR, or the structural-delta metrics. Per model (FDR across the model$\times$metric family), the NASR effect is significant for all four base models but varies widely, from $\eta^2 = 0.29$ for gemma to $0.02$ for Qwen. As the null model below shows, in our setup even this small association does not reflect behavioral coupling.

\paragraph{Is the Topology Effect Behavioral or Mechanical?}
A neighborhood-structured metric such as NASR can move with topology for two very different reasons: agents' answers come to track those of the neighbors the graph gives them (behavioral), or the graph structure and marginal answer counts mechanically constrain the metric regardless of who neighbors whom. To separate these, we run a response-shuffling null model. For each E3 run we hold the graph fixed and permute which agent occupies which node, applying the \emph{same} permutation to all 11 snapshots, so each agent keeps its full opinion trajectory while its spatial relationship to its neighbors is destroyed. We recompute every neighborhood-structured metric under the real assignment and under 200 such permutations, and report where the real value falls in its own null distribution.

The result calls for caution when interpreting neighborhood metrics. Across all 336 E3 runs the real value of every neighborhood-structured metric sits almost exactly on its own null. Only $6.5\%$ of runs have a real NASR outside its permutation null at $p < 0.05$, barely above the $5\%$ expected by chance, with a mean $z$ of $-0.06$ and a mean real value ($0.078$) indistinguishable from the null ($0.079$).
% MOVED TO APPENDIX: the same test on the other neighborhood-structured metrics. The other metrics behave the same way: mean local agreement ($7.4\%$ significant, mean $z \approx 0.00$), cross-cutting edge fraction ($4.8\%$, below chance), and assortativity ($7.4\%$, mean $z = 0.20$).
Crucially, the response-shuffled null reproduces the entire per-topology NASR ordering, from fully connected and cycle at the top to power-law cluster and Barab\'asi--Albert at the bottom, with null means within $0.001$ of the real means at every topology (Table~\ref{tab:null}).

\begin{table}[t]
\centering
\small
\setlength{\tabcolsep}{4pt}
\begin{tabular}{lrrr}
\toprule
Topology & NASR real & NASR null & mean $z$ \\
\midrule
fully connected & 0.090 & 0.092 & $-0.19$ \\
cycle & 0.090 & 0.090 & $-0.12$ \\
stochastic block & 0.083 & 0.083 & $+0.04$ \\
Erd\H{o}s--R\'enyi & 0.080 & 0.080 & $+0.07$ \\
forest fire & 0.070 & 0.070 & $-0.02$ \\
Barab\'asi--Albert & 0.069 & 0.070 & $-0.02$ \\
power-law cluster & 0.064 & 0.065 & $-0.21$ \\
\bottomrule
\end{tabular}
\caption{Null-testing the impact of topology on NASR. Response-shuffling null model (E3, 48 runs per topology, 200 permutations each). Permuting which agent sits on which node, while preserving each agent's answer trajectory, leaves NASR unchanged (real$-$null $\leq 0.002$). Mean $z$ is the real NASR in standard deviations of its own null, $(\text{NASR}_{\text{real}} - \mu_{\text{null}})/\sigma_{\text{null}}$; values near $0$ mean real and shuffled are indistinguishable.}
\label{tab:null}
\end{table}

In our setup, then, the small topology effect on neighbor-alignment metrics is essentially \emph{mechanical}. A topology fixes the size and structure of each agent's neighborhood, which largely determines what the neighbor-majority metrics can be; the measured differences between topologies are not evidence that agents' answers track those of their graph neighbors. The near-perfect real--null agreement at every topology shows that, here, the effect is about structure rather than behavior.

We are careful about the scope of this conclusion. The null tests whether \emph{the topology effect we measured} requires behavioral coupling; it does not show that topology cannot produce such coupling. Our regime is one in which coupling would be hard to detect: agents are stateless and observe one recent thread per turn, answers are read out through a binary survey, and graphs are density-matched. Topology could plausibly matter more with persistent agent memory, longer horizons, stronger exposure to neighbors, recommender-mediated feeds, dynamic rewiring, or outcomes such as information diffusion rather than survey answers. Topology is known to matter in human networks and in classical opinion-dynamics models, and nothing here contradicts that. What our result does establish is a practical point: a raw topology effect on a neighborhood metric is not, on its own, evidence of social influence, and a response-shuffling null is a cheap way to check. Section~\ref{sec:crosssim} repeats this check in three other simulators.

%\paragraph{Density.} Within the five density-matched parameterizable graph families, the empirical mean degree is only weakly associated with NASR (Pearson $r = 0.17$, $p = 0.010$, $n = 240$) and not significantly associated with the other metrics. Combined with the null model, this locates the small measured effect in the mechanical structure of neighborhoods rather than in either edge count or genuine behavioral coupling. \emph{Flag: implementation-tied and largely mechanical; the direction happens to match classical results but is not, in this simulator, behaviorally driven.}


\subsection{Survey Context and the News-Agent Dose}
\label{sec:context}

Adding an agent's own prior survey answers to its context (survey-in-context) is a mid-sized driver: in the E1-only pass it explains $\eta^2_{\text{partial}} = 0.55$ of the order-consistency rate and $0.18$ of OSR. Mechanistically, agents that see the survey question tend to make every message about it, which sharpens response dynamics.

The news dose (0, 1, or 4 broadcasters) is significant but small: for OSR $\eta^2 = 0.06$ ($q = 0.02$), and smaller still for MFR and NASR. Adding more broadcasters nudges population-level metrics but does not qualitatively reshape the dynamics, consistent with a persistent broadcaster acting orthogonally to local neighbor structure. \emph{Flag: implementation-tied (context mechanic; broadcaster placed at the highest-degree node).}

\subsection{Homophily}
\label{sec:homophily}

Homophilic initialization is the dominant term for the change in assortativity over a run (E1-only: partial $\eta^2 = 0.28$ on $\Delta$ assortativity, $q \approx 10^{-30}$), as expected since assortativity measures homophily. But how that clustering persists or erodes is model-gated: model$\times$homophily ($\eta^2 = 0.16$) and question$\times$homophily ($\eta^2 = 0.07$) are both significant. Homophily reliably creates structure; how much survives the run depends on the base model, another non-additive interaction. \emph{Flag: direction expected to generalize; magnitude model-dependent.}

\subsection{Proportions: Average versus Distribution}
\label{sec:proportions}

The population-composition scheme matters, but no single scheme dominates. In the E1-only pass, proportions has a small significant main effect on NCC ($\eta^2_{\text{partial}} = 0.064$, $q = 2.5\times10^{-5}$) but a larger model$\times$proportions interaction ($\eta^2 = 0.144$, $q \approx 10^{-10}$), and on OSR the interaction is larger still ($\eta^2 = 0.317$, $q \approx 10^{-30}$). Table~\ref{tab:proportions} shows why: the sign of the consensus effect flips by base model. The \emph{average} scheme is the most consensus-eroding for Qwen ($-0.115$) but the most consensus-\emph{building} for gemma ($+0.110$); \emph{distribution} builds consensus for the Llama models but erodes it for Qwen. There is no column-wise winner.

\begin{table}[t]
\centering
\small
\setlength{\tabcolsep}{3.5pt}
\begin{tabular}{lrrrr}
\toprule
Proportions & Llama-8B & Minitaur & Qwen & gemma \\
\midrule
average & $-0.055$ & $-0.042$ & $-0.115$ & $+0.110$ \\
blueprint & $-0.051$ & $-0.016$ & $+0.004$ & $+0.132$ \\
distribution & $+0.039$ & $+0.043$ & $-0.060$ & $-0.001$ \\
uniform & $+0.088$ & $+0.026$ & $-0.037$ & $+0.124$ \\
\bottomrule
\end{tabular}
\caption{Mean Net Consensus Change (NCC) by proportions scheme and base model (E1, 36 runs per cell). The effect of population composition on consensus flips sign across models: matching the full human distribution rather than the modal answer matters, but only in interaction with the base model.}
\label{tab:proportions}
\end{table}

\emph{Flag: interaction expected to generalize; specific weights implementation-tied.}

\subsection{Measurement Validity: How Much Answer Change Is Socially Driven?}
\label{sec:measurement}

A change in a stateless agent's survey answer can come from the social content it read, from any other change in its context, or from the phrasing of the question itself. We separate these sources in three ways.

\paragraph{Prompt Sensitivity (Dual-Order Surveys).} Every survey is asked in both option orders. The order-consistency rate, the fraction of answers that survive reversal, is low and strongly model-dependent: Qwen is most consistent ($0.81$ base, $0.87$ fine-tuned), while gemma is least ($0.37$ base, rising to $0.63$ fine-tuned), with the Llama models near $0.41$--$0.50$. In the worst cell more than half of all answers flip on phrasing alone. Order-consistency is itself dominated by model and fine-tuning (model $\eta^2 = 0.39$, question $0.28$, context $0.09$), which means prompt sensitivity is not uniform noise but a systematic, design-dependent property.

\paragraph{No-stimulus Baseline.} With no external stimulus and survey-context off, isolated agents under greedy decoding have OSR exactly $0.000$ for all four models, a sanity check the simulator passes. With survey-context \emph{on}, however, isolated agents drift \emph{more} than social runs (Minitaur reaches OSR $0.51$ in isolation versus $0.15$ socially): seeing one's own prior answer is a strong non-social source of drift, important to control for in context-on runs.

\paragraph{Scrambled-Stimulus Baseline.} With scrambled stimulus, agents observe socially meaningless threads drawn from an unrelated corpus. This still produces OSR of roughly $0.05$--$0.19$; the socially driven excess (normal minus scrambled) is small ($0.04$--$0.08$) and survives FDR correction only for some model$\times$context cells. Only the excess over the scrambled floor should be read as socially driven; much of the raw shift rate is context-perturbation noise, not social influence.

\paragraph{Margin-Stratified Flips.}  The log-probability margin of an answer is the gap in log-probability the model assigns to its chosen option over the next-best one, a measure of how confident that answer is. Flip probability decreases monotonically with the earlier answer's margin: under normal stimulus, answers in the lowest-margin bin flip $32\%$ of the time versus $5.7\%$ in the highest. Low-margin answers behave like noise and high-margin ones like context-driven change, so survey-response metrics should be stratified or attenuated by margin.

Together these baselines justify the terminology of Section~\ref{sec:metrics}: what the survey metrics capture is a mixture of context-driven response change and response-level variability, neither of which is belief revision, and the two should be kept distinct. \emph{Flag: absolute noise-floor magnitudes implementation-tied; the phenomenon (prompt sensitivity of stateless agents) expected to generalize.}


\section{Do the Findings Transfer Across Simulators?}
\label{sec:crosssim}

Everything above was measured in one simulator with one fine-tuning setup, so some findings may describe our implementation rather than LLM social simulation in general. To tell the two apart, we re-ran the central tests in three independently built simulators: OASIS \citep{yang2025oasisopenagentsocial}, where tool-calling agents act on a Twitter-like platform with a recommender and a relational back end; Concordia \citep{vezhnevets2023generative}, whose component-based generative agents keep an associative memory (we use its minimal entity prefab, without a game master); and SiliSocS \citep{sarangi2026ease}, whose EASE-configured agents act through tool calls on a Twitter-like back end with follower-chronological timelines. Each simulator keeps its own prompting, memory, feed construction and action space. Everything else is shared: two of our four base models (Qwen2.5-7B-Instruct and Llama-3.1-Minitaur-8B) with and without the same BluePrint adapters (assigned uniformly, served with vLLM); the same personas, follow graph and activity schedule for a given seed; a neutral trending-topic post shown to every agent at the start; the same scrambled-stimulus corpus; and the same dual-order log-probability survey as in our simulator, administered out of band on each simulator's own representation of the agent's context and never written into agent memory.

The design crosses simulator (3) $\times$ base model (2) $\times$ fine-tuning (on/off) $\times$ question (2, Q2 and Q3) $\times$ seed (2). Each cell runs an Erd\H{o}s--R\'enyi network under normal and scrambled stimulus, plus Barab\'asi--Albert and cycle networks under normal stimulus, for 192 runs of 64 agents over 20 rounds (half of the agents act per round; 6 survey snapshots). The scale is smaller than our main study, so we test whether each effect is present and in which direction, not its exact magnitude. Statistics follow Section~\ref{sec:stats}: Type-III partial $\eta^2$ within each simulator, bootstrap CIs on signed differences, and the same response-shuffling null (200 permutations per run).

\begin{table*}[t]
\centering
\small
\setlength{\tabcolsep}{4pt}
% PLACEHOLDER: replaced by crosssim/results/table_xsim_summary.tex (crosssim/slurm/collect.sh copies it here)
\begin{tabular}{lrrrrrrr}
\toprule
 & \multicolumn{3}{c}{Fine-tuning (OSR)} & Scrambled & Order & Flip rate & NASR \\
\cmidrule(lr){2-4}
Simulator & $\eta^2_p$ & $\Delta$ Llama & $\Delta$ Qwen & floor share & consist. & low$\to$high margin & real$-$null (\% sig.) \\
\midrule
Ours (V2, $N{=}256$) & 0.75 & $+$ & $-$ & -- & 0.37--0.87 & 0.32$\to$0.06 & $-0.001$ (6.5) \\
\midrule
OASIS & \xs{--} & \xs{--} & \xs{--} & \xs{--} & \xs{--} & \xs{--} & \xs{--} \\
Concordia & \xs{--} & \xs{--} & \xs{--} & \xs{--} & \xs{--} & \xs{--} & \xs{--} \\
SiliSocS & \xs{--} & \xs{--} & \xs{--} & \xs{--} & \xs{--} & \xs{--} & \xs{--} \\
\bottomrule
\end{tabular}

\caption{Cross-simulator replication (192 runs). Fine-tuning: partial $\eta^2$ of fine-tuning status on OSR within each simulator (model, fine-tuning$\times$model, graph and question in the model) and the signed OSR change from turning the adapters on, per base model. Scrambled floor share: OSR under scrambled stimulus divided by OSR under normal stimulus (Erd\H{o}s--R\'enyi runs). Order consistency: range over model$\times$fine-tuning cells. Flip rate: share of answers that flip by the next snapshot, lowest vs.\ highest margin quintile. NASR: mean real minus response-shuffled null, with the share of runs significant at $p<.05$ in parentheses. The first row repeats the corresponding values from our simulator.}
\label{tab:xsim}
\end{table*}

\paragraph{What transfers.} \xs{[Fill from crosssim/results/report.md. Claim 1: fine-tuning is the largest OSR/NASR term in X of 3 simulators (partial $\eta^2$ ...), and its sign is/is not model-gated (Llama ..., Qwen ...). Claim 2: scrambled-stimulus OSR is ...\% of normal-stimulus OSR across simulators. Claim 3: order consistency ranges from ... to ..., again depending on model and fine-tuning. Claim 4: flip rate falls/does not fall with margin in X of 3 simulators. Claim 5: in every simulator the per-topology NASR differences are reproduced by the response-shuffled null (real$-$null ..., ...\% of runs significant).]}

\paragraph{What does not.} \xs{[List the findings that differ across simulators (e.g.\ magnitude of the fine-tuning effect, direction for one model, size of the scrambled floor), and the simulator mechanics that plausibly explain each (tool-call action spaces, recommender exposure, memory design, raw completion vs.\ chat prompting).]}

We therefore flag findings as \emph{cross-simulator} when they hold in our simulator and in all three external ones, and as \emph{implementation-tied} otherwise; Section~\ref{sec:discussion} uses these labels. All four simulators share stateless LLM agents, short horizons and survey-based measurement, so even cross-simulator findings are scoped to that regime.

\section{Discussion}
\label{sec:discussion}

Three findings organize our results. First, what dominates survey-response dynamics is the agent itself, not the social world around it. Fine-tuning on social-media data is the single largest lever (partial $\eta^2$ up to $0.89$), and the base model matters chiefly through its interactions with question and population, not as a standalone effect. The structural parameters one might expect to be decisive, network topology, homophily, and activity level, move these metrics very little in the ranges we swept.

Second, the space is mostly additive with a few sharp exceptions. Nearly all parameter pairs act independently, and the non-additivity that exists is carried by a small set of interactions (model$\times$question, model$\times$proportions, model$\times$homophily, fine-tuning$\times$population) in which the base model gates the effect of another choice, so the same decision can move a metric in opposite directions on two models. One cannot assume an effect measured on one model transfers to another, nor add marginal effects as if independent.

Third, and most consequential for how the field reports results, much of the apparent answer change in a stateless simulator is not socially driven. Our null models separate socially driven response change from artifacts, and the gap is large: the measured topology effect is reproduced by a response-shuffling null, much of the raw shift rate survives no better than a scrambled-stimulus baseline, and low-confidence answers flip on option order alone. What looks like social influence can be bookkeeping about neighborhood structure or prompt-level noise.

We are careful about what generalizes. The cross-simulator replication (Section~\ref{sec:crosssim}) supports \xs{[fine-tuning dominance, the model-gated direction of its effect, the scrambled-stimulus floor, the prompt sensitivity of stateless agents, margin-dependent flips, and the mechanical character of the measured topology effect --- edit to match results]} across OASIS, Concordia, SiliSocS and our own simulator. Homophily seeding assortativity and the interaction-borne role of model family were not re-tested and remain claims about our simulator. The exact noise-floor magnitudes, survey-context and news-agent mechanics, and the cadence artifact are implementation-tied. All of these findings are scoped to stateless agents, short horizons and survey-based measurement; simulators with persistent agent state or different outcome measures may behave differently. We make no realism claim: our deliverable maps which parameters move outcomes and how reliably, not which region is most human-like.

The practical takeaway is a short checklist. Fix and report the base model and fine-tuning status first, since they dominate. Report an order-consistency rate and a scrambled-stimulus floor alongside any survey-response metric, and stratify by answer confidence. Before attributing a neighborhood metric to topology, run a response-shuffling null and density-match graphs. Match survey cadence to population size when studying scale. None of these is expensive, and together they separate what a simulator does from what its measurement apparatus manufactures.


\section{Limitations}


Model identity bundles parameter scale, tokenizer, and instruction tuning; only the Minitaur-versus-Llama-3.1-8B pair is a clean within-architecture contrast, and even it confounds post-training recipe. Our agents have no persistent state beyond their context window and adapter weights, so ``survey-response dynamics'' is not belief revision: we measure how answers depend on context, and any reading of these results in terms of beliefs or persuasion would need agents with explicit state and a different measurement design. Likewise, our topology conclusions concern the metrics and regime we study; they do not establish that network structure is behaviorally inert in LLM simulations in general. We estimate main effects and two-way interactions only. The full design-space sweep was run in a single simulator; the cross-simulator replication re-tests only its central claims, at a smaller scale ($N{=}64$, 20 rounds), with two of the four base models and two of the three questions, and all four simulators share stateless agents and survey-based measurement. Non-additivity beyond the model$\times$fine-tuning interaction, and every finding about parameters we did not re-test, should be read as properties of our simulator. Finally, we characterize internal sensitivity and make no realism claim, which would require matched human-trajectory data that does not exist at this scale.


\section*{Ethical Statement}

Our work studies controlled simulations of synthetic LLM agents with no individual-level user data. The intended use is methodological: to help researchers build more rigorous simulators and diagnose design artifacts. Insights about which choices amplify response shifts could in principle inform platform manipulation, but the synthetic setting limits direct transfer, and the same understanding is a prerequisite for defensive measures. Our measurement-validity tooling is meant in part to reduce over-claiming of persuasion or belief-change effects the underlying system does not support.


\paragraph{AI usage disclosure}\textit{ This manuscript and its accompanying code were co-written with Claude Code. The authors have reviewed and edited the entire text and take responsibility for all content.}

\bibliography{aaai2027}

\end{document}
